\section{Related Work}

Automatic statistics management predates learned cardinality estimation.
Chaudhuri and Narasayya studied cost-driven automation of the statistics set
needed by a query optimizer \cite{ChaudhuriNarasayya2001}. CORDS discovers
correlations and soft functional dependencies through sampling and uses them
to guide column-group statistics \cite{Ilyas2004CORDS}. Our goal is related,
but the what-if evaluation is performed by PostgreSQL's native planner and the
deployment artifact is explicitly a stock PostgreSQL recommendation.

DB2 LEO closes a feedback loop around optimizer mistakes and learned
information \cite{Stillger2001LEO}. Query-feedback dependency detection is a
nearby line of work: Haas et al. rank attribute dependencies from sparse query
feedback to guide statistics configuration
\cite{Haas2007AttributeDependencies}. The present
study instead freezes an offline workload snapshot and evaluates native
statistics designs without adding an online feedback estimator.

PostgreSQL's extended-statistics documentation describes the native catalog
objects, MCV lists, and functional dependencies that motivate this work
\cite{PostgreSQL16Docs,PostgreSQL16Catalog}. HypoPG demonstrates a practical
PostgreSQL what-if mechanism for hypothetical indexes \cite{HypoPG}; our
patched substrate targets statistics payloads and planner selectivity paths,
not indexes. EDB Query Advisor provides a software precedent for PostgreSQL
statistics recommendations and explicitly explores column combinations
\cite{EDBQueryAdvisor}. These precedents reinforce why a statistics advisor
alone is not the novelty claim here. The research question is whether a
catalogless native-statistics substrate, a production-oriented boundary, and a
sample-to-full-data protocol provide useful empirical leverage.
